\chapter{Experiments}
\label{chap:experiments}

This chapter presents the experimental evaluation of all systems developed in this master's dissertation. First, the shared experiment settings are described in Section \ref{sec:exp_settings}, including the custom dataset split, dataset statistics, and hardware configuration. Then, experiments are presented for the user-side recommendation system (Section \ref{sec:exp_user}), the recruiter-side recommendation system (Section \ref{sec:exp_recruiter}), and the congestion balancing strategies (Section \ref{sec:exp_congestion}).

% ============================================================
\section{Experiment Settings}
\label{sec:exp_settings}
% ============================================================

\subsection{Custom Train/Validation/Test Split}
\label{sec:custom_split}

% Move section 2.2.1 from old chapter 2 here.
% Explain: why window-based split instead of within-window split,
% cold-start motivation, ignoring the original Train/Test user distinction.
% Include figure fig:custom_split.
TODO

\subsection{Dataset Statistics and Characteristics}
\label{sec:dataset_stats}

% Fill in the empty section 2.4 from old chapter 2 here.
% Include: number of users, jobs, applications per window,
% distribution of applications per user, sparsity of the interaction matrix.
TODO

\subsection{Hardware and Software Configuration}
\label{sec:hardware}

% Describe the machine used for training (GPU type, RAM, etc.)
% PyTorch version, Python version, any relevant library versions.
TODO

% ============================================================
\section{User-Side Recommendation Experiments}
\label{sec:exp_user}
% ============================================================

\subsection{Experiment Goal}
\label{sec:exp_user_goal}

% One paragraph: what are you trying to show?
% Goal: build a two-tower model that ranks relevant jobs highly for each user,
% evaluated on Window 7 as a cold-start scenario.
TODO

\subsection{Experiment Setting}
\label{sec:exp_user_setting}

% Describe the specific setup for these experiments:
% - Negative sampling strategy (city-based, 4:1 ratio) and motivation
% - Training users sampled (100,000)
% - Hyperparameters: lr=0.001, batch_size=1024, output_dim=64, patience=5
% - Early stopping criterion (validation AUC)
% - Evaluation restricted to Window 7 jobs at inference
TODO

\subsection{Ablation Study}
\label{sec:exp_user_ablation}

% Move ablation table and discussion from chapter 3 here.
% 2x2 comparison: address vs no address, norm vs no norm.
% Explain findings: normalization hurts, address helps after fix-location.
% Mention the broken encoding finding and its impact.
TODO

\begin{table}[htbp]
    \centering
    \caption{Ablation study results after 3 training epochs (AUC-ROC on validation set).}
    \label{tab:ablation_user}
    \begin{tabular}{|l|l|c|c|c|}
        \hline
        \textbf{Location Features} & \textbf{L2 Normalization} & \textbf{Val Loss} & \textbf{Accuracy} & \textbf{AUC} \\
        \hline
        Excluded & No  & 0.3513 & 84.60\% & 0.8591 \\
        Excluded & Yes & 0.4687 & 84.18\% & 0.7954 \\
        Included & No  & 0.2317 & 91.24\% & \textbf{0.9365} \\
        Included & Yes & 0.4693 & 84.13\% & 0.7950 \\
        \hline
    \end{tabular}
\end{table}

\subsection{Final Model Results}
\label{sec:exp_user_results}

% Report final validation metrics for the winning configuration.
% Val Loss: 0.2233, AUC: 0.9461, stopped at epoch 11.
% Fill in remaining metrics (Accuracy, Hit@K, Recall@K, NDCG@K) once available.
TODO

\begin{table}[htbp]
    \centering
    \caption{Final model validation results (location features included, no L2 normalization).}
    \label{tab:user_final_results}
    \begin{tabular}{|l|c|}
        \hline
        \textbf{Metric} & \textbf{Value} \\
        \hline
        Validation Loss & 0.2233 \\
        Accuracy        & TODO\% \\
        AUC-ROC         & 0.9461 \\
        Hit@5           & TODO \\
        Recall@5        & TODO \\
        NDCG@5          & TODO \\
        Hit@10          & TODO \\
        Recall@10       & TODO \\
        NDCG@10         & TODO \\
        \hline
    \end{tabular}
\end{table}

\subsection{Qualitative Inference Evaluation}
\label{sec:exp_user_inference}

% Add 3-4 example users from Window 7 with their actual applications
% and what the model recommended. Discuss whether recommendations make sense
% qualitatively (job title match, career progression, location).
% Also discuss the effect of same-state-first post-processing.
TODO

% ============================================================
\section{Recruiter-Side Recommendation Experiments}
\label{sec:exp_recruiter}
% ============================================================

\subsection{Experiment Goal}
\label{sec:exp_recruiter_goal}

% One paragraph: what are you trying to show?
% Goal: evaluate whether the lightweight logistic regression model 
% can reliably predict LLM-generated suitability labels,
% and whether the features capture meaningful candidate-job fit signals.
TODO

\subsection{Experiment Setting}
\label{sec:exp_recruiter_setting}

% Describe the specific setup:
% - How many candidate-job pairs were generated
% - LLM used (Qwen 3.5 35B via Ollama), label distribution
% - Negative sampling strategy for balanced dataset
% - JobID-based train/val/test split (80/10/10)
% - 5-fold stratified cross-validation setup
% - Features used (86 total: heuristic + semantic)
TODO

\subsection{Results}
\label{sec:exp_recruiter_results}

% Report cross-validation and test set results.
% Metrics: Accuracy, Precision, Recall, F1, AUC-ROC.
% Discuss feature importance / model coefficients if available.
% Compare Logistic Regression vs Random Forest vs XGBoost vs Ensemble.
TODO

\begin{table}[htbp]
    \centering
    \caption{Recruiter-side model comparison (5-fold cross-validation).}
    \label{tab:recruiter_results}
    \begin{tabular}{|l|c|c|c|c|c|}
        \hline
        \textbf{Model} & \textbf{Accuracy} & \textbf{Precision} & \textbf{Recall} & \textbf{F1} & \textbf{AUC} \\
        \hline
        Logistic Regression & TODO & TODO & TODO & TODO & TODO \\
        Random Forest       & TODO & TODO & TODO & TODO & TODO \\
        XGBoost             & TODO & TODO & TODO & TODO & TODO \\
        Ensemble            & TODO & TODO & TODO & TODO & TODO \\
        \hline
    \end{tabular}
\end{table}

% ============================================================
\section{Congestion Balancing Experiments}
\label{sec:exp_congestion}
% ============================================================

\subsection{Experiment Goal}
\label{sec:exp_congestion_goal}

% One paragraph: what are you trying to show?
% Goal: measure congestion in the current recommendation system
% and evaluate whether the hybrid strategy reduces it for both sides
% without significantly hurting recommendation quality.
TODO

\subsection{Experiment Setting}
\label{sec:exp_congestion_setting}

% Describe the setup:
% - Use Window 7 as the evaluation window
% - Windows 1-6 as historical data for congestion analysis
% - How congestion is measured (job-side: application concentration,
%   user-side: recommendation overlap)
% - Metrics used to evaluate balancing effectiveness
TODO

\subsection{Job-Side Congestion Analysis}
\label{sec:exp_job_congestion}

% Present the job-side congestion analysis results.
% Show distribution of applications per job in Window 7.
% Identify over-exposed jobs.
TODO

\subsection{User-Side Congestion Analysis}
\label{sec:exp_user_congestion}

% Present the user-side congestion analysis results.
% Show recommendation overlap across users.
TODO

\subsection{Hybrid Balancing Results}
\label{sec:exp_hybrid_results}

% Present the results of the combined strategy.
% Show before/after congestion metrics.
% Show impact on recommendation quality (does relevance drop?).
TODO